\documentclass[12pt]{article}
\usepackage{geometry}
\geometry{margin=1in}
\usepackage{hyperref}
\usepackage{listings}
\usepackage{xcolor}
\usepackage{fancyhdr}
\usepackage{graphicx}

\pagestyle{fancy}
\fancyhead[L]{Northstar Studio}
\fancyhead[R]{Project Documentation}
\fancyfoot[C]{\thepage}

\lstset{
  language=bash,
  basicstyle=\ttfamily\small,
  breaklines=true,
  keywordstyle=\color{blue},
  commentstyle=\color{gray},
  stringstyle=\color{red},
  showstringspaces=false,
  tabsize=2,
  frame=single,
  rulecolor=\color{black!30}
}

\title{\textbf{Northstar Studio} \\ Portfolio \& Order Management System \\ \textit{Complete Project Documentation}}
\author{}
\date{\today}

\begin{document}

\maketitle

\tableofcontents
\newpage

\section{Project Overview}

Northstar Studio is a freelance web development portfolio and B2B order management system (OMS). The platform allows clients to:
\begin{itemize}
  \item Browse portfolio projects and case studies
  \item Submit project inquiries with service type, budget, and timeline
  \item Track order status and communicate with the development team
  \item Process payments via M-Pesa integration
\end{itemize}

The internal admin portal enables the development team to:
\begin{itemize}
  \item View all incoming client inquiries and orders
  \item Update order status and add internal notes
  \item Manage portfolio projects and case studies
  \item Process payments and invoicing
\end{itemize}

\section{Technology Stack}

\subsection{Frontend}
\begin{itemize}
  \item \textbf{React 18}: Component-based UI framework
  \item \textbf{Vite}: Fast build tool and dev server
  \item \textbf{React Router v7}: Client-side routing
  \item \textbf{Tailwind CSS}: Utility-first CSS framework
  \item \textbf{Lucide React}: Icon library
\end{itemize}

\subsection{Backend}
\begin{itemize}
  \item \textbf{Node.js}: JavaScript runtime
  \item \textbf{Express.js}: Web application framework
  \item \textbf{MongoDB}: NoSQL document database
  \item \textbf{Mongoose}: MongoDB object modeling
  \item \textbf{JWT}: JSON Web Tokens for authentication
\end{itemize}

\subsection{Security \& Middleware}
\begin{itemize}
  \item \textbf{Helmet.js}: HTTP headers security
  \item \textbf{CORS}: Cross-Origin Resource Sharing protection
  \item \textbf{express-rate-limit}: API rate limiting
  \item \textbf{express-mongo-sanitize}: NoSQL injection protection
  \item \textbf{Morgan}: HTTP request logging
  \item \textbf{express-validator}: Input validation
\end{itemize}

\section{Project Structure}

\begin{verbatim}
my_website/
├── src/                          # Frontend React application
│   ├── api/
│   │   └── client.js            # Axios HTTP client configuration
│   ├── components/
│   │   ├── Layout.jsx           # Main layout wrapper (ADMIN BUTTON REMOVED)
│   │   ├── ProjectCard.jsx      # Portfolio project display card
│   │   ├── SectionHeading.jsx   # Reusable section heading component
│   │   └── SubmissionForm.jsx   # Client inquiry form component
│   ├── context/
│   │   └── AuthContext.jsx      # JWT token & user state management
│   ├── pages/
│   │   ├── HomePage.jsx         # Landing page with portfolio & inquiry form
│   │   ├── AdminLoginPage.jsx   # Admin sign-in page
│   │   └── AdminDashboardPage.jsx # Admin order & project management
│   ├── App.jsx                  # Router configuration
│   ├── main.jsx                 # React entry point
│   └── index.css                # Global styles
├── config/
│   └── db.js                    # MongoDB connection with fallback to in-memory DB
├── controllers/
│   ├── authController.js        # Admin login & profile endpoints
│   ├── orderController.js       # Order submission & management
│   ├── projectController.js     # Portfolio project CRUD operations
│   ├── uploadController.js      # File upload handling
│   └── paymentController.js     # M-Pesa payment processing
├── middleware/
│   ├── authMiddleware.js        # JWT verification & admin role protection
│   ├── errorMiddleware.js       # Centralized error handling
│   └── validation.js            # Express-validator error handler
├── models/
│   ├── User.js                  # Admin user schema with password hashing
│   ├── Order.js                 # Client inquiry & order schema
│   ├── Project.js               # Portfolio project schema
│   └── Payment.js               # Payment transaction schema
├── routes/
│   ├── authRoutes.js            # POST /api/auth/login, GET /api/auth/me
│   ├── orderRoutes.js           # POST /api/orders, GET/PUT/DELETE for admin
│   ├── projectRoutes.js         # Portfolio CRUD endpoints
│   ├── uploadRoutes.js          # File upload endpoint
│   └── paymentRoutes.js         # M-Pesa STK push, callback, verify
├── utils/
│   └── seedAdmin.js             # Initialize default admin account
├── server.js                    # Express app bootstrap & startup
├── .env                         # Environment variables (local)
├── .env.example                 # Environment variables template
├── package.json                 # Dependencies & NPM scripts
├── tailwind.config.js           # Tailwind CSS configuration
├── vite.config.js               # Vite build configuration
└── README.md                    # Quick start guide
\end{verbatim}

\section{Admin Access \& Security}

\subsection{Hidden Admin Interface}

The admin button has been removed from the public-facing UI to prevent unauthorized access attempts. The admin portal is still fully functional but \textbf{completely hidden} from the public interface.

\textbf{To access the admin panel:}
\begin{enumerate}
  \item Navigate to \texttt{http://localhost:5173/admin/login} directly in your browser
  \item Use your email and password credentials
  \item Upon successful authentication, you will be redirected to the admin dashboard
\end{enumerate}

\subsection{Security Measures}

The admin system implements industry-standard security practices:

\subsubsection{Authentication}
\begin{itemize}
  \item \textbf{JWT Tokens}: Admin sessions are protected by JSON Web Tokens stored securely in localStorage
  \item \textbf{Password Hashing}: Admin passwords are hashed using bcrypt with 10 salt rounds
  \item \textbf{Role-Based Access Control}: Only users with the \texttt{admin} role can access protected endpoints
\end{itemize}

\subsubsection{Route Protection}
\begin{itemize}
  \item All admin API endpoints require valid JWT authentication via Bearer token
  \item The \texttt{authMiddleware.js} validates every protected request and verifies the admin role
  \item Invalid or expired tokens are rejected automatically
\end{itemize}

\subsubsection{API Rate Limiting}
\begin{itemize}
  \item General API rate limit: \textbf{100 requests per 15 minutes}
  \item Order submission rate limit: \textbf{5 requests per hour}
  \item Prevents brute force attacks and abuse
\end{itemize}

\subsubsection{Input Validation}
\begin{itemize}
  \item All client-submitted data is validated using express-validator
  \item Email format validation ensures legitimate contact information
  \item Phone numbers and project details are sanitized
  \item NoSQL injection attacks are prevented via express-mongo-sanitize
\end{itemize}

\subsubsection{HTTP Security Headers}
\begin{itemize}
  \item Helmet.js enforces HTTPS, XSS protection, and clickjacking prevention
  \item CORS is strictly configured to allow only authorized origins
  \item Content Security Policy headers protect against script injection
\end{itemize}

\section{Client Submission Process}

\subsection{Overview}

Clients submit project inquiries through the website contact form. The submission process is simple, intuitive, and includes robust validation.

\subsection{Step-by-Step Submission Workflow}

\subsubsection{Step 1: Navigate to Contact Form}
The client scrolls to the ``Contact'' section on the homepage and finds the project inquiry form. The form is prominently displayed and easy to locate.

\subsubsection{Step 2: Fill Out Project Details}
The client provides the following information:
\begin{itemize}
  \item \textbf{Client Name}: Full name (required, non-empty)
  \item \textbf{Company Name}: Business name (optional)
  \item \textbf{Email Address}: Valid email for communication (required, must be valid email format)
  \item \textbf{Phone Number}: Mobile or office phone (required, non-empty)
  \item \textbf{Project Type}: Selection from predefined list
    \begin{itemize}
      \item Business Website
      \item E-Commerce Store
      \item Web App
      \item Landing Page
      \item Full-Stack System
    \end{itemize}
  \item \textbf{Package Plan}: Selection from three tiers
    \begin{itemize}
      \item Starter: KSh 25,000
      \item Growth: KSh 60,000
      \item Custom: From KSh 120,000
    \end{itemize}
  \item \textbf{Budget Range}: Client's budget estimate (required)
  \item \textbf{Timeline}: Expected project duration
    \begin{itemize}
      \item 2-4 Weeks
      \item 4-8 Weeks
      \item 8-12 Weeks
      \item 12+ Weeks
    \end{itemize}
  \item \textbf{Project Details}: Detailed description of project scope and requirements (required, minimum 20 characters)
  \item \textbf{Payment Method}: Currently supports M-Pesa
  \item \textbf{M-Pesa Phone Number}: Mobile number for payment (if M-Pesa selected)
\end{itemize}

\subsubsection{Step 3: Form Validation}
The form performs real-time validation:
\begin{itemize}
  \item All required fields are checked before submission
  \item Email address format is validated
  \item Project details minimum length is enforced (20 characters)
  \item Invalid entries display helpful error messages to the user
\end{itemize}

\subsubsection{Step 4: Submit Inquiry}
The client clicks the ``Submit'' button. The frontend displays a loading state and sends the form data to the backend API endpoint:

\texttt{POST /api/orders}

\subsubsection{Step 5: Backend Processing}
The server receives the inquiry and:
\begin{enumerate}
  \item Validates all input against security rules
  \item Checks rate limits (max 5 submissions per hour from an IP)
  \item Sanitizes the data to prevent injection attacks
  \item Stores the order record in MongoDB with default status: \texttt{pending}
  \item Assigns a unique order ID and timestamp
  \item Returns a success response with order confirmation
\end{enumerate}

\subsubsection{Step 6: Client Confirmation}
The client receives:
\begin{itemize}
  \item On-screen success message confirming receipt
  \item Order ID for reference
  \item Confirmation that the Northstar Studio team will review the inquiry
\end{itemize}

\subsection{Order Status Flow}

After submission, orders pass through the following status stages:

\begin{lstlisting}
pending  →  under_review  →  quoted  →  accepted  →  in_progress  →  completed
                    ↓
            rejected (optional)
\end{lstlisting}

\subsection{Common Troubleshooting}

\subsubsection{Submission Not Going Through}
\begin{itemize}
  \item Verify all required fields are filled (marked with *)
  \item Ensure email address is valid (correct format: user@domain.com)
  \item Check project details meet 20 character minimum
  \item Clear browser cache and reload the page
  \item Try submitting from a different browser
\end{itemize}

\subsubsection{``Too Many Requests'' Error}
\begin{itemize}
  \item Rate limit: maximum 5 submissions per hour per IP address
  \item Solution: Wait at least 1 hour before submitting again, or use a different network
\end{itemize}

\subsubsection{Payment Issues}
\begin{itemize}
  \item Ensure M-Pesa phone number is in valid Kenya format: +254... or 07...
  \item Confirm account has sufficient M-Pesa balance
  \item Check internet connection stability
  \item For Sandbox testing, use test phone numbers provided in .env
\end{itemize}

\section{Admin Dashboard Features}

\subsection{Authentication}

Admin credentials are seeded automatically at server startup:

\begin{lstlisting}
Email:    admin@company.com
Password: ChangeMe123!
\end{lstlisting}

\textbf{Important:} Change this password immediately after first login in a production environment.

\subsection{Dashboard Capabilities}

\subsubsection{Orders Management}
\begin{itemize}
  \item View all submitted client inquiries and orders
  \item Filter by status (pending, under\_review, quoted, accepted, in\_progress, completed)
  \item Update order status to communicate progress
  \item Add internal notes and remarks for team reference
  \item Delete irrelevant or spam orders
  \item Export order data for analytics
\end{itemize}

\subsubsection{Projects Management}
\begin{itemize}
  \item Create, edit, and delete portfolio projects
  \item Add project title, slug, description, and technical stack
  \item Upload featured image and project gallery
  \item Add links to live website and GitHub repository
  \item Mark projects as featured for homepage prominence
  \item Manage project categories for filtering
\end{itemize}

\subsubsection{Metrics \& Analytics}
\begin{itemize}
  \item View total orders and revenue metrics
  \item Track order fulfillment rate
  \item Monitor conversion funnel from inquiry to completed project
  \item See average project timeline and budget
\end{itemize}

\section{API Endpoints Reference}

\subsection{Authentication}

\begin{lstlisting}
POST   /api/auth/login              Login with email & password
GET    /api/auth/me                 Get current admin profile
\end{lstlisting}

\subsection{Orders}

\begin{lstlisting}
POST   /api/orders                  Submit new client inquiry
GET    /api/orders                  List all orders (admin only)
PUT    /api/orders/:id/status       Update order status (admin only)
DELETE /api/orders/:id              Delete order (admin only)
\end{lstlisting}

\subsection{Projects}

\begin{lstlisting}
GET    /api/projects                List all portfolio projects
POST   /api/projects                Create new project (admin only)
PUT    /api/projects/:id            Update project (admin only)
DELETE /api/projects/:id            Delete project (admin only)
\end{lstlisting}

\subsection{Payments}

\begin{lstlisting}
POST   /api/payments/stk-push       Initiate M-Pesa payment
POST   /api/payments/callback       M-Pesa payment callback
GET    /api/payments/verify/:id     Verify payment status
\end{lstlisting}

\section{Environment Variables}

Create a \texttt{.env} file in the project root with:

\begin{lstlisting}
# Server Configuration
PORT=5000
NODE_ENV=development

# Database (optional - uses in-memory fallback if not available)
MONGODB_URI=mongodb://localhost:27017/northstar-studio

# JWT Secret (change in production)
JWT_SECRET=your_jwt_secret_key_here

# CORS Origins
ALLOWED_ORIGINS=http://localhost:5173,http://localhost:3000

# M-Pesa (Daraja) Configuration
MPESA_CONSUMER_KEY=your_consumer_key
MPESA_CONSUMER_SECRET=your_consumer_secret
MPESA_SHORTCODE=your_shortcode
MPESA_PASSKEY=your_passkey
MPESA_CALLBACK_URL=https://yourdomain.com/api/payments/callback

# File Upload
MAX_FILE_SIZE=5242880
ALLOWED_FILE_TYPES=image/jpeg,image/png,image/webp

# Email Service (for notifications)
SMTP_HOST=smtp.gmail.com
SMTP_PORT=587
SMTP_USER=your_email@gmail.com
SMTP_PASS=your_app_password
\end{lstlisting}

\section{Installation \& Deployment}

\subsection{Local Development Setup}

\begin{lstlisting}[language=bash]
# Clone or extract project
cd my_website

# Install dependencies
npm install

# Start backend server (runs on port 5000)
npm run server

# In a new terminal, start frontend dev server (runs on port 5173)
npm run client

# Or run both concurrently (if available)
npm run dev
\end{lstlisting}

\subsection{Production Build}

\begin{lstlisting}[language=bash]
# Create optimized production build
npm run build

# Output is in dist/ directory - ready for deployment
\end{lstlisting}

\subsection{Database Setup}

\begin{itemize}
  \item \textbf{Development}: In-memory MongoDB (mongodb-memory-server) - no setup needed
  \item \textbf{Production}: Configure MONGODB\_URI to point to hosted MongoDB instance
    \begin{itemize}
      \item MongoDB Atlas (cloud) - recommended
      \item Self-hosted MongoDB server
    \end{itemize}
\end{itemize}

\section{Security Best Practices}

\subsection{Before Going to Production}

\begin{enumerate}
  \item Change the default admin password immediately
  \item Generate a secure, random JWT\_SECRET (minimum 32 characters)
  \item Enable HTTPS/TLS on your domain
  \item Configure ALLOWED\_ORIGINS to only include your domain
  \item Set NODE\_ENV=production
  \item Enable MongoDB authentication with strong database password
  \item Set up automated backups for database
  \item Configure rate limits appropriately for expected traffic
  \item Use environment variables for all secrets (never commit to git)
  \item Set up CORS whitelist to prevent cross-site attacks
  \item Implement request logging and monitoring
\end{enumerate}

\subsection{Ongoing Security}

\begin{itemize}
  \item Regularly update Node.js packages: \texttt{npm audit fix}
  \item Monitor admin login attempts for suspicious activity
  \item Implement IP whitelisting for admin access if possible
  \item Use strong, unique passwords for all admin accounts
  \item Rotate JWT secrets periodically
  \item Keep database backups and test restoration regularly
  \item Review and audit order data for suspicious submissions
  \item Monitor API usage and rate limit violations
\end{itemize}

\section{Support \& Maintenance}

For technical issues or feature requests:
\begin{enumerate}
  \item Check the troubleshooting section
  \item Review server logs for error messages
  \item Verify all environment variables are configured correctly
  \item Ensure MongoDB connection is working
  \item Test API endpoints with Postman or similar tools
\end{enumerate}

\vfill

\begin{center}
\textit{Last Updated: \today}\\
Northstar Studio - Freelance Web Development
\end{center}

\end{document}
